% 第12章（本部分第5章）图：Arnoldi 关系 A V_k = V_{k+1} \bar H_k 的矩阵形状（示意取 d=7, k=4）。
% \bar H_k 为 (k+1)×k 上 Hessenberg：h_{i,j} 仅在 i<=j+1 处可能非零。
% 高亮第 j=2 列：A v_2 = h_{12} v_1 + h_{22} v_2 + h_{32} v_3。
\begin{tikzpicture}[x=1mm,y=1mm,
  note/.style={font=\FigureNoteFont\fontsize{7.5}{9}\selectfont,text=FigureMuted}]
\def\c{3}\def\d{7}\def\k{4}
% A：d×d
\draw[FigureStroke,fill=FigureBlueFill,line width=.6pt] (0,0) rectangle (21,-21);
% V_k：d×k，第2列高亮
\begin{scope}[shift={(24,0)}]
\foreach \j in {1,...,4}{
  \ifnum\j=2 \def\fc{FigureGreenFill}\else\def\fc{white}\fi
  \draw[draw=FigureGrid,fill=\fc,line width=.4pt] ({(\j-1)*\c},0) rectangle ({\j*\c},-21);
}
\draw[FigureStroke,line width=.6pt] (0,0) rectangle (12,-21);
\node[note] at (4.5,2.4) {$\bm v_2$};
\end{scope}
\node at (42,-10.5) {$=$};
% V_{k+1}：d×(k+1)，前三列与第2列的系数对应
\begin{scope}[shift={(46,0)}]
\foreach \j in {1,...,5}{
  \ifnum\j<4 \def\fc{FigureGreenFill}\else\def\fc{white}\fi
  \draw[draw=FigureGrid,fill=\fc,line width=.4pt] ({(\j-1)*\c},0) rectangle ({\j*\c},-21);
}
\draw[FigureStroke,line width=.6pt] (0,0) rectangle (15,-21);
\node[note] at (4.5,2.4) {$\bm v_1,\bm v_2,\bm v_3$};
\end{scope}
% \bar H_k：(k+1)×k
\begin{scope}[shift={(64,0)}]
\foreach \i in {1,...,5}{\foreach \j in {1,...,4}{
  \pgfmathtruncatemacro{\nz}{\i<=\j+1 ? 1 : 0}
  \ifnum\nz=1
    \ifnum\j=2 \def\fc{FigureGreen}\else\def\fc{FigureBlueFill}\fi
  \else\def\fc{white}\fi
  \draw[draw=FigureGrid,fill=\fc,line width=.4pt] ({(\j-1)*\c},{-(\i-1)*\c}) rectangle ({\j*\c},{-\i*\c});
}}
\draw[FigureStroke,line width=.6pt] (0,0) rectangle (12,-15);
\node[note] at (4.5,2.4) {$h_{1,2},h_{2,2},h_{3,2}$};
\draw[FigureMuted,line width=.4pt] (12.8,-13.5) -- (15,-13.5) node[note,anchor=west] {$h_{k+1,k}$};
\end{scope}
% 下方标注
\foreach \x/\n/\s in {10.5/{\bm A}/{d\times d},30/{\bm V_k}/{d\times k},53.5/{\bm V_{k+1}}/{d\times(k+1)}}{
  \node at (\x,-25) {$\n$};
  \node[note] at (\x,-29.5) {$\s$};
}
\node at (70,-19) {$\overline{\bm H}_k$};
\node[note] at (70,-23.5) {$(k+1)\times k$};
\node[note] at (70,-27.5) {上Hessenberg};
\end{tikzpicture}
